% ============================================================
% OP_Nc_anomaly_classification_v1.tex
% GUT1.1 rewritten per GPT review: "N_c-dependent anomaly
% classification and no-overclaim lemma". Own derivation;
% literature page-verified. NO preprint edits by this file.
% ============================================================
\section*{$N_c$-dependent anomaly classification (no-overclaim
lemma); GUT1.1 v2}

\subsection*{0. Reframing per review}
Goal is NOT ``anomalies pin $N_c=3$''.
Goal: classify what anomaly cancellation does and does not fix
when $N_c$ varies, and state explicitly which inherited inputs are
needed to make $N_c=3$ physical.
Possible (and, as shown below, actual) outcome: anomaly
constraints do \emph{not} select $N_c=3$ without observed-charge
or rank input --- a negative result that still improves the
article by preventing overclaim.

\subsection*{1. Own derivation: the anomaly system at general
$N_c$}
Content per generation (left-handed Weyl):
$Q\in(\mathbf N_c,\mathbf 2)_{Y_Q}$,
$u^c,d^c\in(\bar{\mathbf N}_c,\mathbf 1)_{Y_{u},Y_d}$,
$L\in(\mathbf 1,\mathbf 2)_{Y_L}$,
$e^c\in(\mathbf 1,\mathbf 1)_{Y_e}$; no $\nu^c$ (the $B{-}L$ flat
direction with free-$Y$ singlet carries over verbatim from AN3).
Conditions:
\begin{align*}
[SU(N_c)]^3&:\ 2-1-1=0\quad\text{(vector-like; any }N_c\text{)},\\
[SU(N_c)]^2U(1)&:\ 2Y_Q+Y_u+Y_d=0,\\
[SU(2)]^2U(1)&:\ N_cY_Q+Y_L=0,\\
\mathrm{grav}^2U(1)&:\ 2N_cY_Q+N_c(Y_u{+}Y_d)+2Y_L+Y_e=0
 \ \Rightarrow\ Y_e=2N_cY_Q,\\
[U(1)]^3&:\ \text{with }Y_{u}=-Y_Q+\delta,\ Y_d=-Y_Q-\delta:\quad
 6N_c\,Y_Q\,(N_c^2Y_Q^2-\delta^2)=0 .
\end{align*}
Hence for \emph{every} $N_c\ge2$, apart from the all-zero and the
degenerate $Y_Q=0$ branches (structure as in AN3), there is an
SM-like family:
\[
Y_L=-N_cY_Q,\quad Y_e=2N_cY_Q,\quad
\{Y_u,Y_d\}=\{-(N_c{+}1)Y_Q,\ (N_c{-}1)Y_Q\},
\]
reducing at $N_c=3$, $Y_Q=\tfrac16$ to eq.~hypercharges
(and $6N_cY_Q(N_c^2Y_Q^2-\delta^2)$ to the AN3 form
$18Y_Q(9Y_Q^2-\delta^2)$).
\emph{Witten:} $N_c+1$ left-handed $SU(2)$ doublets per generation
$\Rightarrow$ $N_c$ \textbf{odd} --- not $N_c=3$.
\emph{Electric charges} (normalisation $Q_e=-1$, i.e.\
$Y_Q=1/(2N_c)$):
\[
Q_u=\frac{N_c+1}{2N_c},\qquad Q_d=-\frac{N_c-1}{2N_c}.
\]
\emph{Baryons:} an $N_c$-quark colour singlet with $k$ up-type
quarks has $Q_B=k-\tfrac{N_c-1}{2}$, integer iff $N_c$ odd; a
proton analogue ($Q_B={+}1$) exists for every odd $N_c$ at
$k=\tfrac{N_c+1}{2}$.
So charge integrality and atom neutrality also do \emph{not}
select $N_c=3$.
\emph{Anomalous $\pi^0$ coefficient:}
$N_c(Q_u^2-Q_d^2)=N_c\,(Q_u{-}Q_d)(Q_u{+}Q_d)
=N_c\cdot1\cdot\tfrac1{N_c}=1$ --- $N_c$-independent.

\subsection*{2. Classification (the deliverable)}
\textbf{Case (A): fixed observed charges as inherited input.}
Inserting $Q_u=\tfrac23$ (equivalently $Q_d=-\tfrac13$) by hand
gives $\tfrac{N_c+1}{2N_c}=\tfrac23\Rightarrow N_c=3$ ---
correct but the input does all the selecting work; must be
labelled as inherited observed-charge input.
\textbf{Case (B): $N_c$-dependent anomaly-free assignments.}
With the $N_c$-dependent charges above, the full anomaly system
(triangles + grav + Witten) is satisfied for every odd $N_c$;
charge integrality, neutral atoms, and the $\pi^0\to\gamma\gamma$
coefficient are all $N_c$-blind.
(Literature nuance, page-verified: for $N_f\ge3$ the
Wess--Zumino--Witten sector does expose $N_c$, e.g.\
$\eta\to\pi^+\pi^-\gamma\propto N_c^2$ --- so ``invisible $N_c$''
holds for the two-flavour pion--photon sector, not universally.)
\textbf{Case (C): what can identify or impose $N_c=3$.}
Beyond anomaly cancellation, identifying or imposing $N_c=3$ requires
at least one additional input, such as:
(i)~inherited observed-charge normalisation (case~A);
(ii)~$N_f\ge3$ flavour-anomaly phenomenology
($\eta\to\pi^+\pi^-\gamma$, $R$-ratio class);
(iii)~an ECT-native rank-selection mechanism --- the rank-3
zero-mode programme of OP-GUT1 (eq.~p7\_zero\_mode\_target).

\textbf{No-overclaim lemma (proposed wording).}
\emph{Anomaly cancellation (including the Witten condition)
constrains the $N_c$-dependent charge assignments --- $N_c$ odd,
hypercharge pattern fixed up to normalisation and the
$u^c\!\leftrightarrow\!d^c$ relabelling, degenerate branch as in
the $N_c=3$ analysis --- but does not by itself select $N_c=3$;
identifying $N_c=3$ requires additional input, such as inherited
observed-charge assignments, phenomenological information sensitive
to $N_c$ (for example in the $N_f\ge3$ WZW/flavour-anomaly sector
or the $R$-ratio class), or an ECT-native rank-selection mechanism
(zero-mode programme, OP-GUT1).}

\subsection*{3. Verified references (bibliographic/abstract-level)}
\begin{itemize}
  \item O.~B\"ar, U.-J.~Wiese, ``Can one see the number of
    colors?'', Nucl.\ Phys.\ B \textbf{609}, 225--246 (2001),
    hep-ph/0105258. [Abstract verified verbatim: Witten anomaly
    $\Rightarrow N_c$ odd; triangle anomalies determine
    $N_c$-dependent quark charges; $\pi^0\to\gamma\gamma$ not
    $\propto N_c^2$; $\eta\to\pi^+\pi^-\gamma\propto N_c^2$ for
    $N_f\ge3$.]
  \item A.~Abbas, Phys.\ Lett.\ B \textbf{238}, 344 (1990); and
    ``On the number of colours in QCD'', hep-ph/0009242.
    [Verified: $\pi^0\to2\gamma$ makes no statement on $N_c$ with
    colour-dependent charges; $R$-ratio indicated as the better
    evidence.]
  \item Secondary (optional): U.-J.~Wiese lecture notes
    hep-ph/0502008 (states the $(N_c{\pm}1)/2$ baryon composition
    verbatim); arXiv:hep-ph/0101329 (general arbitrary-$N$
    representation analysis).
\end{itemize}
Proposal: cite B\"ar--Wiese + one Abbas reference; decide on the
optional two at review.

\subsection*{4. Implementation sketch (NOT executed; for separate
approval)}
Placement per review answer to Q2 --- dual role:
\begin{enumerate}
  \item New paragraph after the AN3 subsubsection in
    \S anomaly\_architecture:
    ``$N_c$-dependent classification and what it does not select''
    (derivation summary + lemma; candidate Level~A algebraic
    classification under stated inherited content).
  \item One-sentence constraint echo in the colour-completion
    open-questions paragraph: anomaly cancellation alone cannot
    supply the rank selection; the rank-3 input must come from the
    zero-mode programme or inherited charge data.
  \item $N_{\rm colour}=3$ table row reworded (variant per recon
    \S4(b)): hypothesis~(C) retained for the module rank itself;
    comment gains the classification pointer.
  \item OP-GUT1 registry comment: + ``anomaly cancellation shown
    not to select the colour rank by itself
    (\S anomaly\_architecture)''.
  \item Bib: BarWiese2001, Abbas1990 (+ optional) --- verified
    data only.
  \item GUT1.0 in the same round: $G_2$-fact citations
    (Harvey--Lawson 1982 / Baez 2002; to be page-verified) with
    the review wording: standard mathematical facts cited;
    non-adopted deeper candidate; no status upgrade of ECT colour.
  \item GUT1.3 in the same round: split the closure-targets prose
    into ``\textbf{Definition: the colour zero-mode problem}'' +
    ``\textbf{Open problem OP-GUT1}'' paragraphs inside
    \S colour\_completion; no toy model, no existence claim.
  \item GUT1.2 and the $G_2$ citation pass (GUT1.0): DEFERRED to a
    separate mini-proposal after source verification (review
    decision); not part of this round.
\end{enumerate}

\subsection*{5. Status discipline}
Forbidden: ``ECT derives $SU(3)$''; ``anomalies pin $N_c=3$'';
``$N_c=3$ is derived''; ``P7 is now derived''; ``the medium
carries $G_2$ structure''.
Safe: ``$N_c$-dependent anomaly classification'';
``anomaly constraints alone do not select the colour rank without
additional inherited/rank-selection input'';
``conditional $G_2\Rightarrow SU(3)$-structure kernel'';
``zero-mode programme remains open''.

\subsection*{6. Questions for GPT}
(Q1)~Approve the lemma wording in \S2?
(Q2)~Include the $N_f\ge3$ visibility nuance
($\eta\to\pi^+\pi^-\gamma$) as the honest counterpoint --- yes/no?
(Q3)~References: B\"ar--Wiese + which Abbas entry (PLB~238 vs
hep-ph/0009242 vs both)?
(Q4)~After sign-off: implement items 1--7 of \S4 in one round
(with GUT1.2 included or deferred)?
